\section{Work Plan}
\label{sec:plan}

\subsection{Completed}
\begin{itemize}
  \item Literature review of the reference study~\cite{severson2019} and its
  follow-up~\cite{attia2020}.
  \item Conversion of the 2018-02-20 batch from MATLAB to Python.
  \item Study of the dataset's structure, units and charging protocols.
  \item Cleaning (removal of cells without a label) and exploratory
  analysis.
  \item First computation of the $\dq$ feature and its correlation with
  cycle life.
  \item Problem statement finalised (Section~\ref{sec:problem}).
\end{itemize}

\subsection{Phases}
\begin{enumerate}
  \item \textbf{Data (P1).} Download the three batches of~\cite{matr_data}.
  Update the conversion notebooks to current \texttt{h5py} and apply the
  authors' published cleaning steps (removing faulty cells and joining
  cells that continued into the next batch). Save a lightweight cache
  holding the summaries and the \texttt{Qdlin}/\texttt{Tdlin} curves for
  cycles 1--100.
  \item \textbf{Baseline (P2).} Implement feature sets F1--F3, the
  mean and policy-mean baselines, and the elastic-net models; reproduce the
  errors reported in~\cite{severson2019} (RQ1).
  \item \textbf{Physics features (P3).} Implement F4: dQ/dV peak
  height, position and area; resistance against $\sqrt{N}$; temperature rise.
  \item \textbf{Model comparison (P4).} Run all 7 models on all 4 feature
  sets with cross-validated tuning and repeated seeds (RQ2, RQ3). Run the
  early-prediction study over $N$ (RQ4).
  \item \textbf{Generalisation and interpretation (P5).} Evaluate on the
  2018-02-20 batch (RQ5); compute SHAP values and permutation importance.
  \item \textbf{Write-up (P6).} Final report and presentation.
\end{enumerate}

\subsection{Timeline}
\begin{table}[h]
\centering
\caption{Planned timeline (weeks from start of phase P1).}
\label{tab:timeline}
\begin{tabular}{@{}lll@{}}
\toprule
Weeks & Phase & Milestone \\
\midrule
1--2  & P1, P2 & Data cache ready; baseline reproduced \\
3--4  & P3     & Feature set F4 implemented \\
5--7  & P4     & Full results table; error-vs-$N$ plot \\
8     & P5     & Unseen-batch results; importance plots \\
9--10 & P6     & Final report and presentation \\
\bottomrule
\end{tabular}
\end{table}

\noindent\todo{Map weeks to calendar dates agreed with the guide.}

\subsection{Risks and mitigations}
\begin{itemize}
  \item \textbf{Large raw files (several GB per batch).} Convert each batch
  once and work from the lightweight cache.
  \item \textbf{Overfitting with about 124 samples.} Use regularisation,
  nested cross-validation and fixed test sets, and compare against simple
  baselines.
  \item \textbf{Physics features give no improvement.} Report this as a
  finding about which mechanisms do and do not predict lifetime.
\end{itemize}
